\documentclass[12pt,a4paper]{article}
\usepackage[utf8]{inputenc}
\usepackage[T1]{fontenc}
\usepackage[english]{babel}
\usepackage{amsmath, amssymb, amsthm}
\usepackage{graphicx}
\usepackage{hyperref}
\usepackage{geometry}
\usepackage{booktabs}
\usepackage{float}
\usepackage{mathtools}
\usepackage{physics}
\usepackage{cleveref}
\usepackage{enumitem}
\usepackage{dsfont}

\geometry{margin=1in}
\hypersetup{
    colorlinks=true,
    linkcolor=blue,
    citecolor=blue,
    urlcolor=blue,
}


% Theorem-like environments
\theoremstyle{plain}
\newtheorem{theorem}{Theorem}[section]
\newtheorem{proposition}[theorem]{Proposition}
\newtheorem{lemma}[theorem]{Lemma}
\newtheorem{corollary}[theorem]{Corollary}

\theoremstyle{definition}
\newtheorem{definition}[theorem]{Definition}

\theoremstyle{remark}
\newtheorem{remark}[theorem]{Remark}

\title{\textbf{Topological State Preparation via \(\mathbb{Z}/6\mathbb{Z}\) Superselection}\\[0.5em]
\large Optimal Phases and DSP Isomorphism}
\author{José Ignacio Peinador Sala\\
\small Independent Researcher, Valladolid, Spain\\
\small \href{https://orcid.org/0009-0008-1822-3452}{ORCID: 0009-0008-1822-3452}\\
\small \texttt{joseignacio.peinador@gmail.com}}
\date{\today}

\begin{document}
\maketitle

\begin{abstract}
Quantum algorithms for arithmetic problems, such as Shor's factorization, typically initialize registers in a uniform superposition, incurring a large Shannon entropy overhead from computationally sterile trajectories. The modular ring \(\mathbb{Z}/6\mathbb{Z}\) provides a natural superselection rule: all primes \(p>3\) satisfy \(p \equiv 1 \pmod{6}\) or \(p \equiv 5 \pmod{6}\). Restricting a quantum register to these \emph{resonant} channels can reduce the effective search space by two thirds. We investigate the optimal phase modulation required to confine amplitude within the resonant channels without violating unitarity. Using high-precision numerical optimization of the amplitude function \(P(x) \propto \exp[A\sin(2\pi x/6 + \phi)]\), we find that the phase maximizing fidelity for the class \(5 \pmod{6}\) is exactly \(\phi_2 = \pi\) rad, while for the class \(1 \pmod{6}\) it is \(\phi_1 = 0\) rad. We prove that \(\phi_2 = \pi\) is a rigorous consequence of the inversion symmetry \(5 \equiv -1 \pmod{6}\) and the conservation of modular parity, requiring no reference to the Riemann zeta function, holography, or Berry phases. The phase structure is isomorphic to the polyphase decomposition of digital signal processing, guaranteeing perfect spectral isolation. We provide a formal verification of the underlying algebraic properties using the Lean 4 theorem prover. These results establish a rigorous foundation for arithmetic-aware quantum state preparation.
\end{abstract}
\vspace{1em}
\noindent\textbf{Keywords:} Quantum State Preparation, Topological Superselection, Modular Arithmetic (\(\mathbb{Z}/6\mathbb{Z}\)), Phase Modulation, Digital Signal Processing, Polyphase Decomposition, Formal Verification, Lean 4.


\section{Introduction}
Quantum algorithms that exploit arithmetic structure, most notably Shor's algorithm for integer factorization \cite{Shor1994}, rely on an initial uniform superposition over a computational basis of size \(N\). This superposition is generated by applying Hadamard gates to all qubits, an operation of constant depth that maximizes the Shannon entropy of the initial state. While optimal from a circuit complexity perspective, this initialization forces the quantum processor to explore a vast volume of trajectories that are arithmetically irrelevant.

The choice of modulus \(m=6\) is not arbitrary. Beyond its arithmetic role---the smallest modulus that cleanly separates primes \(p>3\) from composite integers---the ring \(\mathbb{Z}/6\mathbb{Z}\) possesses a deep geometric connection: its elements correspond to the sixth roots of unity \(e^{2\pi i r/6}\), which form the vertices of a regular hexagon in the complex plane. This hexagon is the fundamental cell of the \textbf{hexagonal lattice} \(A_2\), the two-dimensional lattice with the highest circle-packing density in the plane \cite{conway1999sphere}. The modular filter thus projects the one-dimensional sequence of computational basis indices onto the optimal packing structure in two dimensions, revealing a geometric order underlying the apparent randomness of discrete phase evaluations. This connection is explored in detail in the companion work on the Modular Spectrum \cite{ModularSpectrum}, where the same polyphase isomorphism is validated in the domain of classical high-precision computation.

A fundamental theorem of elementary number theory states that every prime number \(p > 3\) satisfies \(p \equiv 1 \pmod{6}\) or \(p \equiv 5 \pmod{6}\). The modular ring \(\mathbb{Z}/6\mathbb{Z}\) thus partitions the integers into two \emph{resonant} congruence classes (\(\mathcal{C}_1\) and \(\mathcal{C}_5\)) and four \emph{sterile} classes (\(\mathcal{C}_0, \mathcal{C}_2, \mathcal{C}_3, \mathcal{C}_4\)). In a factorization algorithm, probability amplitude deposited in a sterile class contributes to decoherence and gate overhead without advancing the computation. A topological initialization that confines amplitude to the resonant channels could, in principle, reduce the effective search space by approximately \(66.7\%\).

Recent work has shown that such a confinement can be achieved by modulating the amplitude with a sinusoidal envelope governed by a phase parameter \(\phi\) \cite{Peinador_Isomorfismo_DSP_2025}:
\begin{equation}
P(x) \propto \exp\left[A \sin\left(\frac{2\pi x}{6} + \phi\right)\right] \cdot \mathds{1}_{x \equiv 1,5 \pmod{6}},
\label{eq:amplitude}
\end{equation}
where \(A\) is an effective gain parameter and \(\mathds{1}\) is the indicator function over the resonant channels. For this single-phase model, numerical optimization reveals a clear structure: the phase maximizing the fidelity for class \(\mathcal{C}_5\) converges to exactly \(\pi\) radians.

This article provides a rigorous theoretical foundation for this observation. We prove that \(\phi_2 = \pi\) is a direct consequence of the discrete inversion symmetry of the unit group \((\mathbb{Z}/6\mathbb{Z})^{\times}\) and the conservation of a modular parity operator. The proof is elementary: it uses only the identity \(5 \equiv -1 \pmod{6}\) and the trigonometric identity \(\sin(\theta + \pi) = -\sin\theta\). We further demonstrate that the phase structure is isomorphic to the polyphase decomposition of digital signal processing, providing a practical framework for hardware implementation. Finally, we report a formal machine verification of the underlying algebraic properties using the Lean 4 theorem prover.

The paper is organized as follows. Section~\ref{sec:algebra} reviews the algebraic structure of \(\mathbb{Z}/6\mathbb{Z}\) and classifies the Hilbert space into resonant and sterile channels. Section~\ref{sec:phases} presents the numerical results for the optimal phases and the rigorous derivation of \(\phi_2 = \pi\). Section~\ref{sec:DSP} develops the polyphase isomorphism. Section~\ref{sec:lean} summarizes the Lean 4 formal verification. We conclude in Section~\ref{sec:conclusion}.

\section{Algebraic Structure and Channel Classification}
\label{sec:algebra}

\subsection{The ring \(\mathbb{Z}/6\mathbb{Z}\)}
The set of congruence classes modulo 6,
\begin{equation}
\mathbb{Z}/6\mathbb{Z} = \{0, 1, 2, 3, 4, 5\},
\end{equation}
forms a commutative ring under addition and multiplication modulo 6. The elements coprime to 6, namely \(1\) and \(5\), constitute the \emph{unit group}
\begin{equation}
(\mathbb{Z}/6\mathbb{Z})^{\times} = \{1, 5\} \cong \mathbb{Z}/2\mathbb{Z}.
\label{eq:unitgroup}
\end{equation}
Euler's totient function gives \(\varphi(6) = 2\), confirming that exactly two elements are invertible. The remaining elements \(\{0,2,3,4\}\) are zero divisors: they lack multiplicative inverses.

The fundamental identity governing the dynamics is
\begin{equation}
5 \equiv -1 \pmod{6}.
\label{eq:inverse}
\end{equation}
This is not a convention but an arithmetic necessity: the two generators are additive inverses.

\subsection{Resonant and sterile channels}
A classical theorem \cite{HardyWright} states that every prime number \(p > 3\) satisfies \(p \equiv 1 \pmod{6}\) or \(p \equiv 5 \pmod{6}\). In the quantum register, we map an integer \(x\) to a computational basis state \(\ket{x}\). We define:

\begin{definition}[Resonant and sterile channels]
The \textbf{resonant channels} are the congruence classes \(\mathcal{C}_1\) and \(\mathcal{C}_5\), whose elements are coprime to 6 and can potentially be prime. The \textbf{sterile channels} are \(\mathcal{C}_0, \mathcal{C}_2, \mathcal{C}_3, \mathcal{C}_4\), whose elements are necessarily composite (or zero).
\end{definition}

\begin{theorem}[Prime confinement]
\label{thm:prime_confinement}
Let \(\mathcal{P}_{>3}\) be the set of primes greater than 3. Then \(\mathcal{P}_{>3} \subset \mathcal{C}_1 \cup \mathcal{C}_5\).
\end{theorem}

\begin{proof}
If \(p > 3\) and \(p \not\equiv 1,5 \pmod{6}\), then \(p \equiv 0,2,3,4 \pmod{6}\). In each case, \(p\) would share a non-trivial factor with 6 (2 or 3), contradicting primality.
\end{proof}

Thus, the probability amplitude deposited in sterile channels is computationally wasted: it contributes to the Shannon entropy of the register without advancing the arithmetic search. A topologically optimized initialization should suppress these channels while preserving unitarity and circuit efficiency.

\section{Optimal Phases for Channel Confinement}
\label{sec:phases}

\subsection{Amplitude modulation and numerical optimization}
We consider the amplitude density function \eqref{eq:amplitude}. For a given gain parameter \(A\) (which acts as an effective inverse temperature), the phase \(\phi\) controls how the sinusoidal envelope aligns with the discrete lattice of integers. For each congruence class \(r \in \{1,5\}\), we define the probability fidelity
\begin{equation}
\mathcal{F}_r(\phi) = \frac{\sum_{x \in \mathcal{C}_r} P(x; \phi)}{\sum_{x \in \mathcal{C}_1 \cup \mathcal{C}_5} P(x; \phi)},
\label{eq:fidelity}
\end{equation}
and seek \(\phi_r^{\text{opt}} = \arg\max_\phi \mathcal{F}_r(\phi)\).

High-precision optimization was performed using a grid search followed by gradient descent on a statevector simulation over the first \(5 \times 10^4\) integers, with gain \(A = 5.0\). The results are summarized in Table~\ref{tab:phases}.

\begin{table}[H]
\centering
\caption{Optimal phases from numerical optimization (\(A=5.0\)).}
\label{tab:phases}
\begin{tabular}{lccc}
\toprule
Class & Optimal Phase \(\phi\) [rad] & Physical Origin & Fidelity \\
\midrule
\(\mathcal{C}_1\) & \(0\) & Symmetry (no shift needed) & \(>0.999\) \\
\(\mathcal{C}_5\) & \(\pi \approx 3.141593\) & Inversion symmetry (\(5 \equiv -1\)) & \(>0.999\) \\
\bottomrule
\end{tabular}
\end{table}

The optimization converges unambiguously: \(\phi_1 = 0\) and \(\phi_2 = \pi\), both to within the numerical precision of the simulation. The two resonant channels are therefore governed by a precise phase relation: \(\phi_2 = \phi_1 + \pi\).

\subsection{Rigorous derivation of \(\phi_2 = \pi\)}
The exact value \(\phi_2 = \pi\) is not a numerical accident but a rigorous consequence of the discrete symmetries of the ring.

\begin{theorem}[Optimal phase for \(\mathcal{C}_5\)]
\label{thm:phi2}
Let the amplitude function be given by Eq.~\eqref{eq:amplitude} with an arbitrary real gain \(A \neq 0\). Then the phase maximizing the fidelity \(\mathcal{F}_5(\phi)\) satisfies \(\phi_2 = \phi_1 + \pi\), where \(\phi_1\) is the maximizer of \(\mathcal{F}_1\). In particular, since the system is perfectly symmetric between the two channels, \(\phi_1 = 0\) and consequently \(\phi_2 = \pi\).
\end{theorem}

\begin{proof}
Consider the total partition function over the resonant subspace:
\begin{equation}
Z(\phi) = \sum_{x \in \mathcal{C}_1 \cup \mathcal{C}_5} \exp\left[A \sin\left(\frac{2\pi x}{6} + \phi\right)\right].
\label{eq:Z}
\end{equation}
Under the transformation \(\phi \mapsto \phi + \pi\), the trigonometric identity \(\sin(\theta + \pi) = -\sin(\theta)\) yields
\begin{equation}
\sin\left(\frac{2\pi x}{6} + \phi + \pi\right) = -\sin\left(\frac{2\pi x}{6} + \phi\right).
\label{eq:sin_shift}
\end{equation}
For \(\phi = 0\), the bijection \(x \leftrightarrow 6-x\) pairs each term in \(\mathcal{C}_1\) with a term in \(\mathcal{C}_5\) of equal magnitude but opposite sign in the sine argument:
\[
\sin\left(\frac{2\pi(6-x)}{6}\right) = \sin\left(2\pi - \frac{2\pi x}{6}\right) = -\sin\left(\frac{2\pi x}{6}\right).
\]
The shift by \(\pi\) exactly compensates this sign flip, giving
\begin{equation}
\sum_{x \in \mathcal{C}_5} \exp\left[A \sin\left(\frac{2\pi x}{6} + \pi\right)\right] = \sum_{x \in \mathcal{C}_1} \exp\left[A \sin\left(\frac{2\pi x}{6}\right)\right],
\label{eq:pairing}
\end{equation}
and vice versa. Consequently, the total partition function is invariant:
\begin{equation}
Z(\phi + \pi) = Z(\phi).
\label{eq:Zperiod}
\end{equation}
The conditional probabilities therefore satisfy the exact duality
\begin{equation}
\mathcal{F}_1(\phi) = \mathcal{F}_5(\phi + \pi), \qquad \mathcal{F}_5(\phi) = \mathcal{F}_1(\phi + \pi).
\label{eq:duality}
\end{equation}
If \(\phi_1\) maximizes \(\mathcal{F}_1\), then \(\mathcal{F}_5(\phi_1 + \pi) = \mathcal{F}_1(\phi_1)\) is the maximal attainable value for \(\mathcal{F}_5\). Hence \(\phi_2 = \phi_1 + \pi\). By direct evaluation, \(\phi_1 = 0\) maximizes \(\mathcal{F}_1\), giving \(\phi_2 = \pi\).
\end{proof}

This proof uses only the inversion \(5 \equiv -1 \pmod{6}\) and the trigonometric identity \(\sin(\theta + \pi) = -\sin\theta\). No reference to the Riemann zeta function, holography, or Berry phases is required. The result is a direct consequence of the \(\mathbb{Z}_2\) symmetry of the unit group.

\subsection{Modular parity operator}
The symmetry underlying Theorem~\ref{thm:phi2} can be expressed in operator form. Define the unitary modular parity operator \(\hat{P}\) acting on the chiral degree of freedom:
\begin{equation}
\hat{P} \ket{k, r} = \ket{k, 6-r},
\label{eq:P}
\end{equation}
where \(x = 6k + r\) with \(r \in \{1,5\}\). Then \(\hat{P}^2 = \mathbb{I}\) and \(\hat{P}\) commutes with any Hamiltonian that treats the two resonant channels symmetrically. The expectation value \(\langle \hat{P} \rangle\) is a conserved quantity, the \emph{modular parity invariant}, ensuring that phase modulation does not induce leakage between the channels.

\subsection{Physical interpretation}
The phase \(\phi_2 = \pi\) is the quantum mechanical manifestation of the additive inversion \(5 \equiv -1 \pmod{6}\). In the \(\mathbb{Z}/6\mathbb{Z}\) ring, the two resonant channels are not independent but are parity partners. The Hadamard-uniform initialization treats them as independent, losing this correlation. The optimal phase restores it, aligning the sinusoidal envelope so that the amplitude in \(\mathcal{C}_5\) is the parity image of the amplitude in \(\mathcal{C}_1\).

\section{Polyphase Isomorphism with Digital Signal Processing}
\label{sec:DSP}

The phase structure described in Section~\ref{sec:phases} has a precise mathematical counterpart in classical signal processing: the polyphase decomposition of discrete-time filters. This section formalizes this connection as an exact mathematical isomorphism, not merely a conceptual analogy.

\subsection{Polyphase decomposition}
Consider a discrete signal \(c_n\) defined for integer \(n \geq 0\). Its \(Z\)-transform is \(X(z) = \sum_{n=0}^{\infty} c_n z^{-n}\). Given a decimation factor \(M\), the signal can be decomposed into \(M\) \emph{polyphase components}:
\begin{equation}
E_r(z) = \sum_{k=0}^{\infty} c_{Mk + r} \, z^{-k}, \qquad r = 0, \dots, M-1.
\label{eq:polyphase}
\end{equation}
The original signal is perfectly reconstructed via
\begin{equation}
X(z) = \sum_{r=0}^{M-1} z^{-r} E_r(z^M).
\label{eq:reconstruction}
\end{equation}
This decomposition is the foundation of multirate filter banks and is widely used in sub-band coding, transmultiplexers, and efficient digital filter implementations \cite{vaidyanathan}.

\subsection{Isomorphism with quantum state preparation}
\begin{theorem}[Polyphase Superselection Isomorphism]
\label{thm:polyphase}
Let the computational basis states \(\ket{x}\) of a quantum register be indexed by integers \(x \geq 0\), and let the amplitude function be \(P(x) = \exp[A\sin(2\pi x/6 + \phi)] \cdot \mathds{1}_{x \equiv 1,5 \pmod{6}}\). The orthogonal projection of the amplitude onto the congruence classes modulo \(M=6\),
\begin{equation}
\ket{\psi_r} = \sum_{k=0}^{\infty} c_{6k+r} \ket{6k+r}, \qquad c_{n} = \sqrt{P(n)},
\label{eq:modes}
\end{equation}
is analytically isomorphic to the polyphase decomposition of the discrete signal \(c_n\) in the \(Z\)-domain.
\end{theorem}

\begin{proof}[Proof sketch]
Define the discrete signal \(c_n = \sqrt{P(n)}\) for \(n \geq 0\). Its \(Z\)-transform is \(X(z) = \sum_{n=0}^{\infty} c_n z^{-n}\). Applying the polyphase decomposition with \(M=6\), the \(r\)-th component is
\[
E_r(z) = \sum_{k=0}^{\infty} c_{6k+r} z^{-k}.
\]
Evaluating at \(z=1\) (which corresponds to summing all amplitudes within the channel), we obtain
\[
E_r(1) = \sum_{k=0}^{\infty} c_{6k+r} = \sum_{x \in \mathcal{C}_r} \sqrt{P(x)}.
\]
The vector of channel amplitudes \((E_0(1), \dots, E_5(1))\) is exactly the vector of modular partial sums of the amplitude function. The reconstruction formula \eqref{eq:reconstruction} guarantees that
\[
X(1) = \sum_{r=0}^{5} E_r(1) = \sum_{r=0}^{5} \sum_{x \in \mathcal{C}_r} \sqrt{P(x)} = \sum_{x \in \mathcal{C}_1 \cup \mathcal{C}_5} \sqrt{P(x)},
\]
with no leakage between channels. The orthogonality of the channels in the index domain (disjoint congruence classes) guarantees that this is a \textbf{perfect reconstruction} system: information deposited in one channel cannot contaminate another.
\end{proof}

\subsection{Implications for quantum hardware}
The mapping to the quantum domain is summarized in Table~\ref{tab:isomorphism}.

\begin{table}[H]
\centering
\caption{Polyphase isomorphism between \(\mathbb{Z}/6\mathbb{Z}\) superselection and DSP.}
\label{tab:isomorphism}
\begin{tabular}{ll}
\toprule
\textbf{Modular Substrate \(\mathbb{Z}/6\mathbb{Z}\)} & \textbf{Digital Signal Processing (\(M=6\))} \\
\midrule
Computational basis index \(\ket{x}\) & Discrete time \(n\) \\
Congruence class \(r = x \bmod 6\) & Polyphase component \(r\) \\
Resonant channels (\(r=1,5\)) & Signal sub-bands (passband) \\
Sterile channels (\(r=0,2,3,4\)) & Suppressed sub-bands (stopband) \\
Phase \(\phi\) & Fractional delay in frequency domain \\
Phase shift \(\Delta\phi = \pi\) & Sign inversion of conjugate sub-band \\
Perfect reconstruction (unitarity) & Perfect reconstruction filter bank \\
\bottomrule
\end{tabular}
\end{table}

This isomorphism guarantees three properties essential for NISQ hardware:

\begin{enumerate}[label=(\roman*)]
    \item \textbf{Unitary isolation:} The Hilbert subspaces associated with different congruence classes are orthogonal. Amplitude can be confined to the resonant channels without violating the unitarity of the active register, exactly as a perfect-reconstruction filter bank preserves signal energy.
    \item \textbf{Local decoupling:} Each polyphase component can be prepared independently, enabling low-depth circuit implementation. This is analogous to the \textit{embarrassingly parallel} nature of polyphase filtering, where each sub-band can be processed on dedicated hardware without inter-channel communication \cite{ModularSpectrum}.
    \item \textbf{Perfect reconstruction:} The phase relation \(\phi_2 = \pi\) (with \(\phi_1 = 0\)) ensures that the conjugate channels combine without destructive interference, preserving the total probability amplitude.
\end{enumerate}

\section{Formal Verification in Lean 4}
\label{sec:lean}

To provide absolute certainty about the algebraic foundations, the core properties of the $\mathbb{Z}/6\mathbb{Z}$ ring have been mechanically verified using the Lean 4 theorem prover with the Mathlib4 library \cite{moura2021lean4, mathlib2020}. Three structural pillars were certified:

\begin{enumerate}
    \item \textbf{Modular involution:} The congruence class $5 \pmod{6}$ is proven to be its own multiplicative inverse ($5 \times 5 \equiv 1 \pmod{6}$) and the additive inverse of 1 ($1 + 5 \equiv 0 \pmod{6}$). These properties are captured by the theorems \texttt{modular\_involution\_mul} and \texttt{modular\_involution\_add}, both proved by direct normalization of modular arithmetic in Lean.
    \item \textbf{Unit group isomorphism:} It is exhaustively certified that the only elements coprime to 6 in $\mathbb{Z}/6\mathbb{Z}$ are $\{1, 5\}$, validating $(\mathbb{Z}/6\mathbb{Z})^{\times} \cong \mathbb{Z}/2\mathbb{Z}$. The proof uses the \texttt{native\_decide} tactic to check all residues $x \in \{0,\dots,5\}$ for the condition $\gcd(x,6)=1$.
    \item \textbf{Topological closure:} The deterministic transition rule $\delta(r, b) = (2r + b) \bmod 6$, which governs the automaton underlying the MPS representation, is proven to be strictly bounded within the ring for all inputs $r \in \{0,\dots,5\}$ and $b \in \{0,1\}$. The theorem \texttt{automaton\_closure} establishes that $(2r+b) \bmod 6 < 6$, and \texttt{automaton\_closure\_full} confirms membership in $\{0,\dots,5\}$.
\end{enumerate}

These formal proofs eliminate any possibility of hidden arithmetic exceptions and provide a certified foundation for the topological state preparation scheme. The full Lean 4 source code is available in the supplementary repository.

\section{Discussion and Conclusion}
\label{sec:conclusion}

We have presented a rigorous analysis of topological state preparation based on the modular ring \(\mathbb{Z}/6\mathbb{Z}\). The optimal phases that confine quantum amplitude within the resonant channels \(\mathcal{C}_1\) and \(\mathcal{C}_5\) have been determined both numerically and analytically: \(\phi_1 = 0\) for class 1, and \(\phi_2 = \pi\) for class 5. The latter is not a numerical fit but a rigorous consequence of the inversion symmetry \(5 \equiv -1 \pmod{6}\) and the \(\mathbb{Z}_2\) structure of the unit group.

The polyphase isomorphism with digital signal processing provides a practical framework for implementing the scheme on quantum hardware: each polyphase component can be prepared independently, and the phase relation \(\phi_2 = \phi_1 + \pi\) guarantees perfect reconstruction. The Lean 4 formal verification certifies the algebraic foundations to the highest standard of rigor.

The result \(\phi_2 = \pi\) is remarkable in its simplicity: it requires only elementary modular arithmetic and basic trigonometry, yet it provides a direct bridge between the algebraic structure of prime numbers and the phase control of quantum registers. This work establishes a foundation for arithmetic-aware quantum state preparation. Obvious next steps include: (a) a complete circuit-level implementation of the polyphase initialization and an analysis of its gate complexity; (b) an investigation of extensions of this superselection scheme to other modular rings relevant to arithmetic algorithms, such as \(\mathbb{Z}/4\mathbb{Z}\) or \(\mathbb{Z}/8\mathbb{Z}\); and (c) the design of adaptive search strategies that exploit the duality \(\mathcal{F}_1(0) = \mathcal{F}_5(\pi)\) to explore both resonant channels with a single circuit.

\section*{Declarations}

\textbf{Funding:} The author declares that no funds, grants, or other support were received during the preparation of this manuscript.

\vspace{0.5em}
\noindent\textbf{Author Contribution:} J.I.P.S. is the sole author of this manuscript. The author independently conceptualized the theoretical framework, performed all mathematical and algebraic derivations, conducted the numerical statevector optimizations, developed the Lean 4 formal verification, and wrote the paper.

\vspace{0.5em}
\noindent\textbf{Ethics Declaration:} Not applicable.

\vspace{0.5em}
\noindent\textbf{Competing Interests:} The author has no relevant financial or non-financial interests to disclose.

\section*{Data Availability}
The numerical simulation data and the Lean 4 source code used for the formal verification of the algebraic properties are available in the supplementary repository associated with this publication on Zenodo. All underlying theoretical derivations and logical steps are explicitly detailed within the manuscript.

\section*{Acknowledgments}
The author expresses deep gratitude to the open-source community for the computational and typesetting tools that facilitated the preparation of this manuscript. Special acknowledgment is given to the developers of the Lean 4 theorem prover and the Mathlib4 library, whose tools enable the rigorous mechanical certification of mathematical foundations. A sincere acknowledgment is also extended to the foundational architects of quantum computation and digital signal processing, whose structural insights provide the unshakeable bedrock for this synthesis. Special thanks are directed to the anonymous peer reviewers and editorial staff, whose critical scrutiny, adversarial rigor, and constructive feedback are essential for maintaining the coherence and integrity of the scientific literature. This work is dedicated to the ongoing advancement of independent, foundational research in quantum information and mathematical physics.

\section*{AI Use Statement}
In accordance with modern scientific transparency standards, artificial intelligence models (specifically DeepSeek and Gemini) were utilized during the preparation of this manuscript. Their application was strictly limited to language editing, structural refinement, \LaTeX\ formatting optimization, and acting as a critical analysis team (\textit{red team}) to challenge, debate, and stress-test the theoretical arguments and logic. All core physical postulates, mathematical deductions, numerical analyses, Lean 4 formalizations, and final conjectures represent the original, independently verified intellectual work of the author.

\section*{Recommended Citation and License}

\noindent\textbf{Recommended Citation:} \\
J.\,I.\,Peinador~Sala, \emph{Topological State Preparation via \(\mathbb{Z}/6\mathbb{Z}\) Superselection: Optimal Phases and DSP Isomorphism}, Zenodo (2026). \url{https://doi.org/10.5281/zenodo.19354010}

\vspace{1em}
\noindent\textbf{Open Access License:} \\
This article is distributed under the terms of the Creative Commons Attribution 4.0 International License (\href{https://creativecommons.org/licenses/by/4.0/}{CC BY 4.0}), which permits unrestricted use, distribution, and reproduction in any medium, provided you give appropriate credit to the original author(s) and the source, provide a link to the Creative Commons license, and indicate if changes were made.

\begin{thebibliography}{99}
\bibitem{Shor1994} P.~W.~Shor, \emph{Algorithms for quantum computation: discrete logarithms and factoring}, Proc.\ 35th Annual Symposium on Foundations of Computer Science (FOCS), IEEE, 124--134 (1994). \url{https://doi.org/10.1109/SFCS.1994.365700}

\bibitem{conway1999sphere} J.~H.~Conway and N.~J.~A.~Sloane, \emph{Sphere Packings, Lattices and Groups}, 3rd ed.\ (Springer-Verlag, 1999).

\bibitem{ModularSpectrum} J.~I.~Peinador~Sala, \emph{The Modular Spectrum of \(\pi\): Theoretical Unification, DSP Isomorphism, and Exascale Validation}, Zenodo (2026). \url{https://doi.org/10.5281/zenodo.17680023}

\bibitem{Peinador_Isomorfismo_DSP_2025} J.~I.~Peinador~Sala, \emph{A Modular DSP Architecture for Extreme-Precision Computation of \(\pi\)}, Zenodo (2026). \url{https://doi.org/10.5281/zenodo.17768718}

\bibitem{HardyWright} G.~H.~Hardy and E.~M.~Wright, \emph{An Introduction to the Theory of Numbers}, 6th ed.\ (Oxford University Press, 2008). \url{https://doi.org/10.1093/oso/9780199219858.001.0001}

\bibitem{vaidyanathan} P.~P.~Vaidyanathan, \emph{Multirate Systems and Filter Banks} (Prentice Hall, 1993).

\bibitem{moura2021lean4} L.~de~Moura and S.~Ullrich, \emph{The Lean 4 theorem prover and programming language}, Proc.\ 28th International Conference on Automated Deduction (CADE 28), Springer, 625--635 (2021). \url{https://doi.org/10.1007/978-3-030-79876-5_37}

\bibitem{mathlib2020} The mathlib Community, \emph{The Lean mathematical library}, Proc.\ 9th ACM SIGPLAN International Conference on Certified Programs and Proofs (CPP 2020), ACM, 367--381 (2020). \url{https://doi.org/10.1145/3372885.3373824}
\end{thebibliography}

\end{document}